\reading{IM-11}{Private Markets Investing}{QUICK WIN}{2}

\srcnote{Cyril Demaria, Maurice Pedergnana, Remy He, Roger Rissi and Sarah
Debrand, \textit{Asset Allocation and Private Markets: A Guide to Investing with
Private Equity, Private Debt, and Private Real Assets} (Wiley, 2021), Chapter 3,
\term{up to section 3.3.10}. No GARP source book covers this reading.}

\begin{imtrapbox}[breakable=false,title={A scope restriction the transcribed spine had dropped}]
GARP restricts this reading to \term{section 3.3.10 and earlier}. That
restriction is in the official learning objective document and was absent from the
transcribed spine this project had been working from. It is restored here.
\end{imtrapbox}

\begin{imkeybox}[title={The six objectives in this reading}]
\begin{enumerate}[label=\textbf{\alph*.}]
\item Explain the motivations for investors to invest in private markets.
\item Describe different investment structures that can be used to invest in private assets and compare their characteristics.
\item Apply and compare different methods to measure the performance of private equity investments and describe their advantages and disadvantages.
\item Explain how different risks can arise for investors in private markets and assess how these risks can be measured.
\item Describe the practice of co-investing in private markets and assess its advantages and disadvantages.
\item Explain benefits and challenges of investing in private markets through an intermediated investment vehicle such as a private equity or private debt fund.
\end{enumerate}
\end{imkeybox}

% ------------------------------------------------------------------
\lo{IM-11 a}{Explain the motivations for investors to invest in private markets.}

Private markets allow private companies to raise capital through equity,
convertible debt and traditional debt. Equity investments in private firms are
\term{private equity (PE)}, debt investments are \term{private debt (PD)}, and
investments in privately held real assets are \term{private real asset (PRA)}
investments. A private equity investment may eventually become publicly traded
through a listing. Some publicly traded firms with limited trading volume raise
capital through private transactions known as \term{private investments in public
equities (PIPEs)}.

\begin{imkeybox}[title={Six motivations}]
\begin{enumerate}
\item \term{Performance and diversification.} Potential for higher returns, plus
diversification from lower correlation with public markets, especially in periods
of uncertainty.
\item \term{Control and value creation.} Investors can actively monitor and
influence the company.
\item \term{Access to specific sectors or stages} unavailable in public markets.
\item \term{Potential for absolute returns.} Returns are often driven by
company-specific improvements and are \emph{less dependent on overall market
movements}.
\item \term{Wider geographic reach.} Access to regions with underdeveloped public
markets, helping diversify globally and \term{reduce home bias}.
\item \term{Specialized investment structures.} Funds of funds and co-investments
allow tailored exposures not typically available in public markets.
\end{enumerate}
\end{imkeybox}

% ------------------------------------------------------------------
\lo{IM-11 b}{Describe different investment structures that can be used to invest in private assets and compare their characteristics.}

An investment structure determines how capital is pooled, managed and governed,
and defines the rights and obligations of investors and managers.

\renewcommand{\arraystretch}{1.2}
\noindent\begin{tabularx}{\textwidth}{@{}>{\raggedright\arraybackslash}p{2.5cm}X@{}}
\toprule
\textbf{\sffamily\small Structure} & \textbf{\sffamily\small Characteristics} \\
\midrule
Direct investment & Investor acquires ownership directly. \term{High control}, but requires significant expertise, active involvement and strong networks. Carries \term{high concentration risk}, limited diversification and information asymmetry. Used by entrepreneurs, family offices and buyout investors; \term{crowdfunding} is also a form of direct investment. \\
Private funds & Capital pooled and managed by \term{general partners} --- the most common structure. Provides diversification but \term{limits individual investor control}. Used by institutional investors and pension funds, typically as \term{closed-end limited partnerships}. \\
Funds of funds & Allocate to a fund that invests in multiple private funds, improving diversification but \term{adding an extra layer of fees}. Used by less experienced investors or those seeking diversification. \\
Securitization & Assets, typically debt, pooled and divided into \term{tranches} with different risk-return profiles. Provides targeted exposure but \term{increases complexity and reduces transparency}. \\
\bottomrule
\end{tabularx}
\renewcommand{\arraystretch}{1}
\figcap{Four structures on one axis: control against diversification. Direct
investment maximises the first and sacrifices the second; funds of funds does the
reverse and pays twice in fees for it. Securitization is the odd one out --- it
does not trade control for diversification but slices risk instead.}

\begin{imkeybox}[title={How the instruments themselves are structured}]
\term{Equity} can be \term{common} or \term{preferred} shares, where preferred may
carry specific rights and protections, leading to different levels of influence.
\term{Debt} can be \term{convertible} or nonconvertible, with different seniority
levels. Convertible debt provides \term{downside protection with upside equity
potential}; senior, subordinate and junior debt define priority of claims.
\end{imkeybox}

% ------------------------------------------------------------------
\lo{IM-11 c}{Apply and compare different methods to measure the performance of private equity investments and describe their advantages and disadvantages.}

\subsection{Internal rate of return}

IRR measures the annualized return on invested capital by incorporating both the
\term{magnitude and timing} of cash flows, making it highly relevant for
long-duration, irregular cash flow investments. It is widely used for comparing
funds and managers.

\begin{imtrapbox}[breakable=false,title={Four ways IRR misleads}]
IRR can be \term{distorted by early distributions}, which artificially increase
the annualized return. It \term{assumes reinvestment at the same IRR}, which is
often unrealistic. Managers may further \term{inflate IRR by delaying capital
calls using credit lines}, effectively shortening the investment period. And it
does not account for risk or scale, focusing only on \emph{deployed} capital.
\end{imtrapbox}

\gapbadge
\srcnote{The source notes cover IRR and stop; the objective's verb is
\emph{compare}, which requires the other methods. Written below from the available
Schweser condensation, one rung below a GARP source.}

\figwrap{%
\begin{tikzpicture}[
  hd/.style={draw=FRMNavy,fill=FRMNavy,rounded corners=1pt,inner sep=4pt,
             font=\scriptsize\sffamily\bfseries,text=white,align=center,
             text width=34mm},
  bx/.style={draw=FRMRule,fill=white,rounded corners=1pt,inner sep=4pt,
             font=\scriptsize\sffamily,align=center,text width=34mm},
  ar/.style={-{Stealth[length=4pt]},FRMGold,line width=0.9pt}]
\node[hd] (irr) at (0,0) {IRR};
\node[bx,below=4mm of irr] (irr1) {annualised; uses \textbf{timing}\\ of cash flows};
\draw[ar] (irr) -- (irr1);
\node[hd] (moic) at (5.0,0) {MOIC};
\node[bx,below=4mm of moic] (moic1) {absolute; \textbf{ignores timing}\\ $1.0\times$ is breakeven};
\draw[ar] (moic) -- (moic1);
\node[bx,below=5mm of moic1] (tvpi) {\textbf{TVPI} $=$ (distributions\\ $+$ residual) / called};
\node[bx,below=2.5mm of tvpi] (dpi) {\textbf{DPI} $=$ distributions / called\\ \emph{cash only, unmanipulable}};
\node[bx,below=2.5mm of dpi] (rvpi) {\textbf{RVPI} $=$ residual / called};
\draw[ar] (moic1) -- (tvpi);
\node[hd] (pme) at (10.0,0) {PME};
\node[bx,below=4mm of pme] (pme1) {benchmarked; cash flows\\ discounted at \textbf{index returns}};
\node[bx,below=2.5mm of pme1] (pme2) {$>1.0$ means the fund\\ beat the benchmark};
\draw[ar] (pme) -- (pme1); \draw[ar] (pme1) -- (pme2);
\end{tikzpicture}}%
{Three families, and what each one adds. IRR is the only measure that uses the
\emph{timing} of cash flows, which is both its strength and every one of its
weaknesses. MOIC discards timing entirely and so cannot be manipulated by
accelerating distributions --- but also cannot be annualised. PME is the only one
that brings in a \emph{benchmark}. The MOIC submetrics form an identity worth
carrying into the exam: \term{TVPI $=$ DPI $+$ RVPI}, so the measure that cannot
be manipulated (DPI, cash only) and the measure that can (RVPI, manager-determined
residual value) always sum to the headline figure.}

\renewcommand{\arraystretch}{1.2}
\noindent\begin{tabularx}{\textwidth}{@{}>{\raggedright\arraybackslash}p{1.5cm}XX@{}}
\toprule
& \textbf{\sffamily\small Advantages} & \textbf{\sffamily\small Disadvantages} \\
\midrule
MOIC & Focuses on absolute performance; independent of dollar size, so comparisons across funds are straightforward; \term{less susceptible to manipulation than IRR}; cash-based submetrics such as DPI are robust. & Measures only \emph{deployed} capital, not committed; not risk-adjusted; \term{ignores time value of money} including holding periods; \term{cannot be annualized}; no explicit benchmark or market context. \\
PME & Incorporates \term{explicit benchmarking} to provide market context; helps neutralise the impact of market cycles; measures manager value creation in a way less prone to manipulation. & Works best with \emph{fully realized} funds; accurate implementation requires matching each distribution with its associated capital call to reflect the correct holding period. \\
\bottomrule
\end{tabularx}
\renewcommand{\arraystretch}{1}
\figcap{Cambridge Associates developed a \term{modified PME (mPME)} to approximate
the distribution-matching calculation; it provides a more robust measure of
manager value creation and is less prone to manipulation.}

\begin{imkeybox}[title={Which measure answers which question}]
\term{Timing matters and you want an annual rate} $\rightarrow$ IRR.
\term{Absolute performance without regard to holding periods} $\rightarrow$ MOIC.
\term{Did the manager beat the market} $\rightarrow$ PME.
\term{Has the money actually come back} $\rightarrow$ DPI.
\end{imkeybox}

% ------------------------------------------------------------------
\lo{IM-11 d}{Explain how different risks can arise for investors in private markets and assess how these risks can be measured.}

\gapbadge
\srcnote{The source notes carry the objective heading and the sub-heading
``Private Market Risks'' on an \emph{otherwise blank page}. The entire objective
is absent. Written below from the available Schweser condensation, one rung below
a GARP source.}

\begin{imgapbox}[title={What the objective asks for}]
Two halves, and the notes supply neither: which risks arise from the
\emph{characteristics} of private assets, and --- the part candidates skip --- what
you actually \emph{measure} to see them.
\end{imgapbox}

\subsection{Where the risks come from}

\renewcommand{\arraystretch}{1.2}
\noindent\begin{tabularx}{\textwidth}{@{}>{\raggedright\arraybackslash}p{2.6cm}X@{}}
\toprule
\textbf{\sffamily\small Risk} & \textbf{\sffamily\small How it arises from the nature of private assets} \\
\midrule
Systematic & Private assets show low correlations with public indices but remain exposed to \term{macroeconomic risks} --- downturns, inflation shifts, interest rate policy. \term{Interest rate risk matters most for leveraged strategies}. Because the assets are illiquid, some market risks such as \term{foreign exchange risk are harder to hedge} than in public markets, and long holding periods complicate hedging further. \\
Illiquidity & Not traded on exchanges, typically \term{locked up for long periods} with limited exit opportunities. \\
Credit & Most relevant for private debt, with \emph{two} components: \term{recovery risk}, the possibility capital is lost if a borrower defaults --- particularly important for distressed investors, though private debt managers have generally limited recovery losses well --- and \term{credit event risk}, where negative developments in public credit markets spill over into private debt. \\
Settlement & Always present when transacting with counterparties, but private market investors conduct \term{extensive due diligence}, which reduces it relative to public market transactions. \\
Information asymmetry & Private firms are not subject to the same regulation as public ones, limiting the amount and quality of available information, which is what makes \term{due diligence} so central. \\
\bottomrule
\end{tabularx}
\renewcommand{\arraystretch}{1}
\figcap{Note the two rows that run against intuition. \term{Settlement risk is
lower} in private markets than in public ones, because of the due diligence the
illiquidity forces. And \term{credit risk has two components}, only one of which
is about the borrower --- the other arrives from the public market the investor
thought they had diversified away from.}

\subsection{How the risks are measured}

\renewcommand{\arraystretch}{1.2}
\noindent\begin{tabularx}{\textwidth}{@{}>{\raggedright\arraybackslash}p{3.0cm}X@{}}
\toprule
\textbf{\sffamily\small Tool} & \textbf{\sffamily\small What it reveals} \\
\midrule
Performance dispersion & \term{Macro risk}: monitor the spread between the highest and lowest pooled average TVPI. Outliers distort it, so a multiyear rolling average helps. A large spread suggests macroeconomic conditions drove returns. Also measures \term{selection risk}, by comparing a fund's TVPI with the pooled average for its category. \\
DPI against TVPI & \term{Reporting risk}: DPI rests entirely on realized cash flows, TVPI on manager-determined NAV estimates. \term{A substantial gap between the two may indicate reporting risk.} \\
Portfolio weights & \term{SAA risk}: compare the investor's weights with a target allocation or a \term{market-neutral portfolio} reflecting the geographic and strategy mix of the relevant allocation. \\
VaR and stress testing & Model worst-case outcomes for strategy-specific risks, including macro shocks and periods of reduced liquidity. \\
Vintage year analysis & The impact of \term{market timing} on performance, and how risks vary across economic cycles. \\
Due diligence reviews & Qualitative: confirm whether the risks that \emph{materialized} match those \emph{expected} at the time of investment. Gaps feed back into future due diligence. \\
Default and loss rates & For private debt, historical default rates and loss statistics benchmark \term{recovery risk}. \\
\bottomrule
\end{tabularx}
\renewcommand{\arraystretch}{1}
\figcap{The measurement half of the objective, which is where the marks are. The
sharpest item is the second: \term{DPI against TVPI is a reporting-risk detector},
because the two differ by exactly the component a manager controls. The first is
the most easily confused --- performance dispersion measures \emph{two} things,
macro risk across the category and selection risk within it, depending on what you
compare against.}

% ------------------------------------------------------------------
\lo{IM-11 e}{Describe the practice of co-investing in private markets and assess its advantages and disadvantages.}

\begin{imdefbox}[title={Co-investing}]
A private investor directly investing in a \term{specific deal} alongside a
private market fund, typically \term{by invitation from a fund manager}. Managers
may extend co-investment opportunities when they need additional capital to
complete a deal that exceeds their usual investment limits. Co-investments often
require \term{quick decision-making and accelerated due diligence}.
\end{imdefbox}

\renewcommand{\arraystretch}{1.2}
\noindent\begin{tabularx}{\textwidth}{@{}XX@{}}
\toprule
\textbf{\sffamily\small Advantages} & \textbf{\sffamily\small Disadvantages} \\
\midrule
\term{Lower fees.} Reduced or no management fees and carried interest, though transaction and due diligence costs may still be shared. & \term{Adverse selection.} Investors may be offered less attractive deals, or deals outside the manager's core expertise, and opportunities may arise when valuations are already high. \\
\term{Greater control.} Select specific deals, controlling portfolio exposures and allocation decisions. & \term{Speed and complexity.} Rapid due diligence and quick decisions, often with limited time for thorough analysis. \\
\term{Higher net returns.} Lower fees often mean higher after-fee returns than fund investments. & \term{Resource and expertise requirements.} Demands strong internal capabilities and the ability to evaluate deals efficiently. \\
\term{Learning and access.} Hands-on experience, and stronger relationships with fund managers. & \term{Concentration risk.} Investments are typically large and concentrated, reducing diversification. \\
& \term{Legal and governance risks.} Limited control over decisions and fewer legal protections than primary fund investors. \\
\bottomrule
\end{tabularx}
\renewcommand{\arraystretch}{1}
\figcap{\term{Adverse selection is the disadvantage to know.} Everything else on
the right-hand column is a cost of doing the work; adverse selection is a
structural problem with the invitation itself. If the manager can place the deal
in the fund, why is it being offered to you?}

% ------------------------------------------------------------------
\lo{IM-11 f}{Explain benefits and challenges of investing in private markets through an intermediated investment vehicle such as a private equity or private debt fund.}

\renewcommand{\arraystretch}{1.2}
\noindent\begin{tabularx}{\textwidth}{@{}XX@{}}
\toprule
\textbf{\sffamily\small Benefits} & \textbf{\sffamily\small Challenges} \\
\midrule
\term{Expertise.} Fund managers bring specialized skills in sourcing, due diligence, structuring, monitoring and exiting investments. & \term{High fees.} Management fees and carried interest reduce net returns and \term{create a performance hurdle}. \\
\term{Risk reduction.} Diversification across deals, sectors and geographies; professional management may reduce losses. & \term{Principal-agent risk.} Misalignment between general partners and limited partners can affect decisions and performance. \\
\term{Deal flow access.} Access to high-quality, often exclusive opportunities unavailable to individual investors. & \term{Lack of transparency.} Limited visibility into underlying investments leads to information asymmetry and reduced control. \\
\term{Cost and resource sharing.} High costs of sourcing, analysis and monitoring are shared across investors. & \term{Blind pool risk.} Capital is committed \emph{without knowing the specific investments}; reliance on manager decisions. \\
& \term{Illiquidity risk.} Capital locked in for long durations; secondary sales may require discounts. \\
& \term{Due diligence demands.} Investors must rigorously evaluate fund strategy, manager track record and governance. \\
& \term{Conflicts of interest.} In capital allocation, co-investments, or holding \emph{both debt and equity} in the same firm. \\
\bottomrule
\end{tabularx}
\renewcommand{\arraystretch}{1}
\figcap{Four benefits against seven challenges --- and the asymmetry is the point.
Every benefit is a service the manager provides; every challenge is a consequence
of having delegated. \term{Blind pool risk} is the one with no analogue in public
markets: you commit the capital before the investments exist. This entire table
sits inside a single image in the source notes and returns nothing to text
extraction.}

\begin{imtrapbox}[title={Consolidated trap summary --- IM-11}]
\begin{itemize}
\item \term{Scope.} GARP examines this reading only \emph{up to section 3.3.10}.
\item \term{Sibling, IM-11 c.} IRR uses timing; MOIC ignores it; PME benchmarks.
Only IRR annualises; only PME brings in the market.
\item \term{Formula, IM-11 c.} \term{TVPI $=$ DPI $+$ RVPI}. DPI is cash and
cannot be manipulated; RVPI is manager-determined residual value and can be.
\item \term{Role, IM-11 c.} Delaying capital calls with credit lines \emph{inflates
IRR} by shortening the measured investment period. It does nothing to MOIC.
\item \term{Polarity, IM-11 d.} Settlement risk is \emph{lower} in private markets
than public ones, because of the due diligence illiquidity forces.
\item \term{Sibling, IM-11 d.} Credit risk has two components: recovery risk from
the borrower, and credit event risk arriving from public markets.
\item \term{Definition, IM-11 d.} A \emph{gap between DPI and TVPI} signals
reporting risk. Performance dispersion signals macro risk and selection risk.
\item \term{Role, IM-11 e.} Co-investments are \emph{by invitation}, which is
precisely why adverse selection is the headline disadvantage.
\item \term{Definition, IM-11 f.} Blind pool risk is committing capital before
knowing the investments --- not illiquidity, and not lack of transparency, which
are listed separately.
\end{itemize}
\end{imtrapbox}

% ==================================================================
